\documentclass[11pt, oneside]{article}
\usepackage{geometry}
\usepackage[T1]{fontenc}
\usepackage[encapsulated]{CJK}
\usepackage{ucs}
\usepackage[utf8]{inputenc}
\usepackage{scalerel}
\usepackage{tikz}
\usetikzlibrary{svg.path}
\usepackage{multicol}
\usepackage{amssymb}


\definecolor{orcidlogocol}{HTML}{A6CE39}
\definecolor{linkedincol}{HTML}{0e76a8}
\tikzset{
  orcidlogo/.pic={
    \fill[orcidlogocol] svg{M256,128c0,70.7-57.3,128-128,128C57.3,256,0,198.7,0,128C0,57.3,57.3,0,128,0C198.7,0,256,57.3,256,128z};
    \fill[white] svg{M86.3,186.2H70.9V79.1h15.4v48.4V186.2z}
                 svg{M108.9,79.1h41.6c39.6,0,57,28.3,57,53.6c0,27.5-21.5,53.6-56.8,53.6h-41.8V79.1z M124.3,172.4h24.5c34.9,0,42.9-26.5,42.9-39.7c0-21.5-13.7-39.7-43.7-39.7h-23.7V172.4z}
                 svg{M88.7,56.8c0,5.5-4.5,10.1-10.1,10.1c-5.6,0-10.1-4.6-10.1-10.1c0-5.6,4.5-10.1,10.1-10.1C84.2,46.7,88.7,51.3,88.7,56.8z};
  }
}

\newcommand\orcidicon[1]{\href{https://orcid.org/#1}{\mbox{\scalerel*{
\begin{tikzpicture}[yscale=-1,transform shape]
\pic{orcidlogo};
\end{tikzpicture}
}{|}}}}

\usepackage{xcolor}
\usepackage{hyperref}
\usepackage{fontawesome}
\hypersetup{
    colorlinks=true,
    linkcolor=blue,
    filecolor=magenta,      
    urlcolor=blue,
}
\definecolor{Mycolor2}{HTML}{820000}
\definecolor{orcidlogocol}{HTML}{A6CE39}
\definecolor{scholar}{HTML}{4285F4}
% 820000
% EFA91D
% EF7C7C
%002147
\renewcommand\labelitemii{\textcolor{Mycolor2}{$\circ$}}
\renewcommand\labelitemiii{\textcolor{Mycolor2}{$\circ$}}
\renewcommand\labelitemi{\textcolor{Mycolor2}{$\circ$}}


\pagestyle{empty}
\geometry{letterpaper,tmargin=0.6in,bmargin=0.8in,lmargin=0.4in,rmargin=0.4in,headheight=0in,headsep=0in,footskip=.3in}

\setlength{\parindent}{0in}
\setlength{\parskip}{0in}
\setlength{\itemsep}{0in}
\setlength{\topsep}{0in}
\setlength{\tabcolsep}{0in}

%%%%%%%%%%%%%%%%%%%%%%%%%%%%%%%%%%%%%%%%%%%%%%%%%%%%%%%%%
% New commands and environments

% This defines how the name looks
\newcommand{\bigname}[1]{
	%\fontfamily{phv}
	%\begin{center}\scshape\Huge#1\end{center}
    {\hspace{50pt}\scshape\Huge#1}
}

% A ressection is a main section (<H1>Section</H1>)
\newenvironment{ressection}[1]{
	\vspace{4pt}
	%\fontfamily{phv}
	{\selectfont\Large\color{Mycolor2}{\textbf{#1}}}
	\vspace{-6.5pt} \\
	\textcolor{Mycolor2}{\rule{\textwidth}{0.7pt}}
	\begin{itemize}
	\vspace{3pt}
}{
	\end{itemize}
}

% A resitem is a simple list element in a ressection (first level)
\newcommand{\resitem}[1]{
	\vspace{-4pt}
	\item[] \begin{flushleft} #1 \end{flushleft}
    %\item[]  {#1} 
}

% A ressubitem is a simple list element in anything but a ressection (second level)
\newcommand{\ressubitem}{
	\vspace{-4pt}
	%\item \begin{flushleft} #1 \end{flushleft}
	\item
}

% A resbigitem is a complex list element for stuff like jobs and education:
%  Arg 1: Name of company or university
%  Arg 2: Location
%  Arg 3: Title and/or date range
\newcommand{\resbigitem}[3]{
	\vspace{-5pt}
	\item[]
	\textbf{#1} --- #2 %\\
	\hfill
	\textit{#3}
}

% This is a list that comes with a resbigitem
\newenvironment{ressubsec}[3]{
	\resbigitem{#1}{#2}{#3}
	\vspace{-2pt}
	\begin{itemize}
}{
	\end{itemize}
}

% This is a list that comes with a resbigitem
\newenvironment{ressubsecP}[3]{
	\resbigitem{#1}{#2}{#3}
	\vspace{2pt}

}

% This is a list that comes with a resbigitem
\newenvironment{ressubsecPu}[1]{
	\resitem{#1}
	\vspace{2pt}

}


% This is a simple sublist
\newenvironment{reslist}[1]{
	\resitem{\textbf{#1}}
	\vspace{-5pt}
	\begin{itemize}
}{
	\end{itemize}
}

\begin{document}
\newcommand{\zhs}[1]{\begin{CJK}{UTF8}{gbsn}#1\end{CJK}}

\fontfamily{ppl} \selectfont


% Name with horizontal rule
\centerline{\huge{\textbf{\textcolor{Mycolor2}{Abel Lawrence Peirson}}}}

\small{
	\vspace{5pt}
	\centerline{\faMapMarker \hspace{1pt} New York, NY \hspace{5pt}
	\faEnvelopeO \hspace{1pt} \href{mailto:alpv95@alumni.stanford.edu}{alpv95@alumni.stanford.edu} \hspace{5pt}
	\faGlobe \hspace{1pt} \href{https://alpv95.github.io}{\texttt{alpv95.github.io}}}
% 	\textcolor{linkedincol}{\faGraduationCap} \hspace{1pt}
% 		\href{https://www.linkedin.com/in/abel-lawrence-peirson-v-b96aba94/}{abel-lawrence-peirson-v}}
	\vspace{1.5pt}
	\centerline{\textcolor{scholar}{\faGraduationCap} \hspace{-1pt}
	\href{https://scholar.google.com/citations?user=MxcIoQYAAAAJ}{abel-lawrence-peirson}
	\hspace{5pt} \faGithub \hspace{1pt} \href{https://github.com/alpv95}{\texttt{alpv95}}}
	\vspace{2pt}
	\centerline{\orcidicon{0000-0001-6292-1911}\hspace{1pt} \href{https://orcid.org/0000-0001-6292-1911}{orcid.org/0000-0001-6292-1911}}
}

\vspace{10 pt}

%%%%%%%%%%%%%%%%%%%%%%%%
\begin{ressection}{\textbf{Education}}
	\vspace{2pt}
	\begin{ressubsec}{\normalsize{Stanford University}}{\emph{Ph.D. in Physics}}{2017 - Mar 2023}
		\ressubitem{Thesis: \href{https://github.com/alpv95/thesis}{Relativistic Jets in the Era of X-ray Polarimetry}.}
		\ressubitem{GPA: 4.0 (in required coursework).}
	\end{ressubsec}

	\begin{ressubsec}{\normalsize{University of Oxford}}{\emph{MPhys in Physics}}{2013 - 2017}
		\ressubitem{First Class Honours with Distinction -- top 10\% of graduating class.}
		\ressubitem{College: Christ Church}
	\end{ressubsec}
\end{ressection}
%%%%%%%%%%%%%%%%%%%%%%%

\begin{ressection}{\textbf{Honors \& Awards}}
	\vspace{2pt}
	\begin{ressubsecP}{Stanford Data Science Scholar (\$100k+)}{Stanford, CA}{2021}
	\end{ressubsecP}

	\begin{ressubsecP}{Future Investigator in NASA Earth and Space Science and Technology FINESST (\$160k+)}{Stanford, CA}{2019}
	\end{ressubsecP}

	\begin{ressubsecP}{Hooke \& Roach Prizes: top STEM student, Christ Church}{Oxford, UK}{2017}
	\end{ressubsecP}

	\begin{ressubsecP}{Oxford International Strategy Scholarship (CERN)}{Oxford, UK}{2016}
	\end{ressubsecP}

	\begin{ressubsecP}{Gold, British Physics Olympiad}{London, UK}{2013}
	\end{ressubsecP}

	\vspace{2pt}
	\resitem{\textbf{\emph{\textcolor{Mycolor2}{Featured in:}}} \href{https://science.nasa.gov/blogs/ixpe/2022/07/11/meet-ixpe-scientist-abel-lawrence-peirson/}{NASA Blog} (2022) --- \href{https://hoggresearch.blogspot.com/2022/11/parameters-and-nuisance-parameters.html}{Hogg's Research} (2022) --- \href{https://podcasts.apple.com/us/podcast/this-astrophysics-grad-student-doesnt-always-make-memes/id1186480811?i=1000419233302}{NVIDIA AI Podcast}, ep.~68 (2018) --- \href{https://techcrunch.com/2018/06/15/dank-learning-system-autogenerates-memes/}{TechCrunch} (2018)}

\end{ressection}

%%%%%%%%%%%%%%%%%%%%%%%%

\begin{ressection}{\textbf{Experience}}

\vspace{2pt}
\begin{ressubsec}{Citadel GQS}{New York / Miami}{2023 - Present  \\ Quantitative Researcher}
	\ressubitem{Alpha research: sequence modelling and structured deep learning.}
\end{ressubsec}

\vspace{2pt}
\begin{ressubsec}{Kavli Institute for Particle Astrophysics and Cosmology}{Stanford University, CA}{July 2018 - Mar 2023 \\ with \textbf{Prof. Roger Romani}}
\ressubitem{State-of-the-art X-ray polarization recovery using simulation-based inference.}
\ressubitem{$>30\%$ improvement in NASA IXPE polarization sensitivity (code adopted as official data analysis pipeline).}
\ressubitem{Quadratic program + nested sampling approach to fit gravitational microlenses.}
		\ressubitem{Testable (and fittable) emission models to explain observed polarization in relativistic plasma jets.}
		%  \ressubitem{Member of the IXPE collaboration: theory, simulation and AGN teams.}
	\end{ressubsec}

\vspace{2pt}
\begin{ressubsec}{Google Brain}{Mountain View, CA}{June - Dec 2022  \\ Student Researcher}
	\ressubitem{Optimization for deep learning, NTK-NNPK correspondence.}
	\ressubitem{Layerwise Bregman exponential families for preconditioning SGD with local Fisher approximations.}
	\ressubitem{MLCommons algorithmic efficiency benchmarking.}
	\ressubitem{Mentors: Rohan Anil / Ehsan Amid / Prof. Manfred K. Warmuth.}
\end{ressubsec}

\vspace{2pt}
\begin{ressubsec}{G-Research}{London, UK}{June - Sep 2021  \\ Quantitative Researcher Intern}
		\ressubitem{Algorithmic trading detection.}
	\end{ressubsec}

\vspace{2pt}
\begin{ressubsec}{Peirson \& Freedman}{Stanford, CA}{Aug 2018 - 2019  \\ Co-founder}
\ressubitem{Built iOS app \textbf{\emph{Dank Learning}} that uses RNNs to generate memes. [\href{https://danklearning.com/}{Webpage}]}
	\end{ressubsec}

\vspace{2pt}
\begin{ressubsec}{Wu Tsai Neurosciences Institute}{Stanford University, CA}{Mar 2018 - Jun 2018 \\ with \textbf{Prof. Shaul Druckmann}}
\ressubitem{Biologically inspired recurrent neural network to reproduce 2D path integration in the Drosophila brain.}
		%  Developing a ring attractor dynamical system neural network to reproduce 2D path integration and pointing vector persistence in the Drosophila brain.} \ressubitem{Results understood using topological data analysis.}
	\end{ressubsec}

\vspace{2pt}
\begin{ressubsec}{NeuroAI Lab}{Stanford University, CA}{Dec 2017 - Mar 2018 \\ with \textbf{Prof. Dan Yamins}}
	\ressubitem{Unsupervised learning: shear-transformation pretraining does not improve image classification.}
		%  Investigating unsupervised visual learning models, namely whether learning affine transformations can be transferred to improved accuracies in object classification, using RotNet and ImageNet datasets.}
	\end{ressubsec}

\vspace{2pt}
\begin{ressubsec}{Department of Physics}{University of Oxford, UK}{Sep 2016 - May 2017 \\ with \textbf{Prof. Garret Cotter}}
% \ressubitem{Modelled blazar spectra and investigated how they change travelling across different intergalactic media.}
		\ressubitem{Placed limits on whether the Cherenkov Telescope Array will constrain the existence of axion-like particles.}
	\end{ressubsec}

\vspace{2pt}
\begin{ressubsec}{CLIC Test Facility}{CERN, Switzerland}{June - Aug 2016 \\ with \textbf{Prof. Philip Burrows}}
\ressubitem{Improved CTF3's Quadrupole scan and reduced uncertainty in the beam energy spread.}
		%  which measures the emittance and Twiss parameters, 
		%  \ressubitem{Reduced the uncertainty in the beam energy spread.} 
% \ressubitem{Designed and implemented new fitting program to enhance beam analysis pipeline.}
		%  Quantified first order chromatic aberration effects. Results were used to sharpen beam injection.}
	\end{ressubsec}

% \vspace{2pt}

% \begin{ressubsec}{Plasma and Fusion Laboratory}{University of Science and Technology of China, Hefei}{June - Aug 2015\\ with \textbf{Prof. Xuan Sun}}
% 	\ressubitem{Improved plasma confinement in the KMAX axisymmetric tandem mirror machine by applying a bias voltage.}
% \end{ressubsec}

\end{ressection}

%%%%%%%%%%%%%%%%%%%%%%%%%%%%%

\begin{ressection}{\textbf{Presentations}}
	\resitem{\textbf{\textcolor{Mycolor2}{--- Invited ---}}}
	\vspace{2pt}

	\begin{ressubsecPu}{\textbf{\emph{IXPE: Science so far}}, \emph{IAUS 375: The Multimessenger Chakra of Blazar Jets}, Dec 2022}
	\end{ressubsecPu}

	\begin{ressubsecPu}{\textbf{\emph{X-ray Polarimetry: A data intensive science}}, \emph{Stanford Data Science Seminar}, 2022}
	\end{ressubsecPu}

	\begin{ressubsecPu}{\textbf{\emph{Tests of gravitational milli-lensing in the blazar PKS 1413+135}}, \emph{Max Planck Institute for Radio Astronomy}, 2022}
	\end{ressubsecPu}

	\begin{ressubsecPu}{\textbf{\emph{Optimal Signal Extraction for IXPE and an Application to Blazars}}, \emph{Naval Research Laboratory Colloquium}, 2021}
	\end{ressubsecPu}

	\begin{ressubsecPu}{\textbf{\emph{Towards Optimal Signal Extraction for IXPE}}, \emph{Third Science Collaboration Meeting (SCM03)}, 2021}
	\end{ressubsecPu}

	\begin{ressubsecPu}{\textbf{\emph{The Polarization Behavior of Synchrotron Self-Compton Emission in Blazars}}, \emph{Understanding the Multiwavelength Blazar Variability - Workshop, Stanford}, 2019}
	\end{ressubsecPu}

	\begin{ressubsecPu}{\textbf{\emph{AI in Design}}, \emph{Used Future: Symposium by Current Obsession}, Pratt Institute NY, 2018}
	\end{ressubsecPu}

	\vspace{2pt}
	\resitem{\textbf{\textcolor{Mycolor2}{--- Contributed ---}}}
	\vspace{2pt}

	\begin{ressubsecPu}{\textbf{\emph{Relativistic Jets in the Era of X-ray Polarimetry}}, \emph{Columbia THEA Seminar}, Nov 2022}
	\end{ressubsecPu}

	\begin{ressubsecPu}{\textbf{\emph{Fishy: Local Fisher Approximation via Shampoo}}, \emph{Google Brain Weekly}, Nov 2022}
	\end{ressubsecPu}

	\begin{ressubsecPu}{\textbf{\emph{Blazar rotations with IXPE}}, \emph{KIPAC Tea}, Oct 2022}
	\end{ressubsecPu}

	\begin{ressubsecPu}{\textbf{\emph{A Geometric Model for Blazar Rotations}}, \emph{Meeting of the European Astronomical Society}, 2020}
	\end{ressubsecPu}

\end{ressection}

%%%%%%%%%%%%%%%%%%%%%%%%%%%%%
%%%%%%%%%%%%%%%%%%%%%%%%

\begin{ressection}{\textbf{Telescope \& Computing Allocations}}
	\vspace{2pt}
	\begin{ressubsecP}{ALMA}{\textnormal{\emph{High-frequency radio observations of IXPE targets}, 7-14 day (Co-I)}}{2021}
	\end{ressubsecP}

	\begin{ressubsecP}{Nordic Optical Telescope}{\textnormal{\emph{Multi-band polarization observations of IXPE targets}, 87ks (Co-I)}}{2021}
	\end{ressubsecP}

	\begin{ressubsecP}{NuSTAR}{\textnormal{\emph{Multi-Energy X-ray observations of IXPE blazar targets}, 140ks (Co-I)}}{2021}
	\end{ressubsecP}

	\begin{ressubsecP}{Swift}{\textnormal{\emph{Monitoring IXPE blazar targets with Swift}, 240ks (Co-I)}}{2020}
	\end{ressubsecP}

	\begin{ressubsecP}{Google Cloud Platform}{\textnormal{\emph{Parametric Density Estimation with Uncertainty using Deep Ensembles}, \$1000}}{2020}
	\end{ressubsecP}

	\begin{ressubsecP}{XMM-Newton}{\textnormal{\emph{Exploring the Synchro-Compton transition in CGRaBS J0211+1051}, 57ks (Co-I)}}{2019}
	\end{ressubsecP}

\end{ressection}

%%%%%%%%%%%%%%%%%%%%%%%%
\begin{ressection}{\textbf{Open Source} ($\bigstar$1000+)}

	\resitem{\textbf{\textcolor{Mycolor2}{Python --- C/C++ --- JAX --- PyTorch --- TensorFlow}}}

	\vspace{4pt}
	\resitem{\textbf{MLCommons Algorithmic Efficiency:} Benchmark and competition for neural network training. [\href{https://github.com/mlcommons/algorithmic-efficiency}{Code}][\href{https://mlcommons.org/en/groups/research-algorithms/}{Webpage}]}

    \vspace{2pt}
	\resitem{\textbf{ixpeobssim:} Simulation and analysis framework for the Imaging X-ray Polarimetry Explorer. [\href{https://github.com/lucabaldini/ixpeobssim}{Code}][\href{https://ixpeobssim.readthedocs.io/en/latest/index.html}{Webpage}]}

    \vspace{2pt}
	\resitem{\textbf{MulensModel:}  Python package for modelling gravitational microlensing events. [\href{https://github.com/rpoleski/MulensModel}{Code}][\href{https://arxiv.org/abs/2201.01110}{Paper}][\href{https://rpoleski.github.io/MulensModel/}{Webpage}]}

	\vspace{2pt}
	\resitem{\textbf{SSCpol:}  Polarized relativistic jet simulation in C with Python wrapper. [\href{https://github.com/alpv95/SSCpol}{Code}][\href{https://ui.adsabs.harvard.edu/abs/2019ApJ...885...76P/abstract}{Paper}]}

	\vspace{2pt}
	\resitem{\textbf{Dank Learning:} `Show and Tell' image captioning for meme generation in TensorFlow. [\href{https://github.com/alpv95/Dank-Learning}{Code}][\href{https://arxiv.org/abs/1806.04510}{Paper}][\href{https://danklearning.com/}{Webpage}]}

\end{ressection}

\newpage

\begin{ressection}{\textbf{Outreach and Service}}

	\begin{ressubsec}{Wonderfest}{Bay Area, CA}{June 2019 - June 2020  \\ Science Envoy}
	\ressubitem{Selected as one of 10 graduate students from Stanford and Berkeley to communicate science to public audiences.}
	\end{ressubsec}

	\begin{ressubsec}{Stanford Diversity and First Generation Office}{Stanford, CA}{Oct 2018 - June 2020  \\ First Generation and Low Income Student Mentor}
		\ressubitem{Mentoring undergraduates from underprivileged backgrounds in all matters of student life.}
	\end{ressubsec}

	\begin{ressubsec}{Stanford Astronomical Society}{Stanford, CA}{Oct 2018 - Oct 2019  \\ Organizing Member}
		\ressubitem{Organizing and leading stargazing events for local schools and the public.}
	\end{ressubsec}

	\begin{ressubsec}{Stanford Science Pen Pals}{Stanford, CA}{Oct 2018 - Oct 2019  \\ Pen Pal}
		\ressubitem{Corresponding about science with 5th graders from a low-income school in Catalonia.}
	\end{ressubsec}

	% 	\begin{ressubsec}{Stanford Judo Club}{Stanford, CA}{June 2019 - June 2020  \\ President}
	%                 \ressubitem{Led the club as captain in competitions.}
	%                 \ressubitem{Set the annual budget and ran all financial matters for both the adults' and children's club.}
	% 	\end{ressubsec}

\end{ressection}

%%%%%%%%%%%%%%%%%%%%%%%%%%%%%

\begin{ressection}{\textbf{Mentoring}}

	\begin{ressubsecP}{Jack Dinsmore (graduate)}{Polarization leakage in the IXPE gas pixel detector}{2022 - Mar 2023}
	\end{ressubsecP}

	\begin{ressubsecP}{Josephine Wong (graduate)}{Simultaneous polarization fitting of the Crab nebula and pulsar}{2021 - Mar 2023}
	\end{ressubsecP}

\end{ressection}

%%%%%%%%%%%%%%%%%%%%%%%
\begin{ressection}{\textbf{Teaching}}
	\vspace{2pt}
	\begin{ressubsecP}{Stanford PHYS113}{\textnormal{\emph{Computational Physics} (Lecturer and Teaching Assistant)}}{Winter 2021}
	\end{ressubsecP}

	\begin{ressubsecP}{Stanford PHYS100}{\textnormal{\emph{Introduction to Observational Astrophysics} (Teaching Assistant)}}{Spring 2019}
	\end{ressubsecP}
\end{ressection}

\begin{ressection}{\textbf{Selected Graduate Coursework}}
	\vspace{-8pt}
	\item[]
	\begin{multicols}{2}
		\begin{itemize}
			\item \textbf{\textit{APPHYS293}} (Theoretical Neuroscience)
			\item \textbf{\textit{CS379C}} (Computational Models of the Neocortex)
			\item \textbf{\textit{CS238}} (Decision Making Under Uncertainty)
			\item \textbf{\textit{STATS207}} (Time Series)
			\item \textbf{\textit{CS224N}} (Natural Language Processing)
			\item \textbf{\textit{CS230}} (Deep Learning)
			\item \textbf{\textit{CME212}} (Advanced Software Development)
			\item \textbf{\textit{EE364a \& b}} (Convex Optimization I \& II)
			\item
			      \textbf{\textit{PHYS266}} (Statistical Methods in Physics)
			\item
			      \textbf{\textit{AA 214}} (Numerical Methods for Compressible Flows)
			\item \textbf{\textit{CS361}} (Engineering Design Optimization)
			\item \textbf{\textit{EE263}} (Linear Dynamical Systems)
		\end{itemize}
	\end{multicols}

\end{ressection}

%%%%%%%%%%%%%%%%%%%

\newpage

\vspace{5pt}

\begin{ressection}{\textbf{Selected Publications} (h-index 44)}

\vspace{4pt}
\resitem{\textbf{\textcolor{Mycolor2}{--- Machine Learning ---}}}
\vspace{2pt}

\begin{ressubsecPu}{\textbf{\emph{Benchmarking Neural Network Training Algorithms
		}} [\href{https://github.com/mlcommons/algorithmic-efficiency}{Code}]\\ G.E.Dahl, et al. incl. \textbf{A.L.Peirson} \href{https://arxiv.org/abs/2306.07179}{2306.07179}, 2023}
\end{ressubsecPu}

\begin{ressubsecPu}{\textbf{\emph{Fishy: Layerwise Fisher Approximation for Higher-Order Neural Network Optimization
			}}\\ \textbf{A.L.Peirson}, E.Amid, Y.Chen, V.Feinberg, M.K.Warmuth, R.Anil \href{https://openreview.net/forum?id=cScb-RrBQC}{\emph{NeurIPS HITY workshop}}, 2022}
\end{ressubsecPu}

\begin{ressubsecPu}{\textbf{\emph{Neural Network Analysis of X-ray Polarimeter Data}}\\ \textbf{A.L.Peirson}, \emph{The Handbook of X-ray and Gamma Ray Astrophysics}, \href{https://link.springer.com/referenceworkentry/10.1007/978-981-16-4544-0_144-1}{\emph{Springer Nature}}, 2022}
\end{ressubsecPu}

\begin{ressubsecPu}{\textbf{\emph{A Deep Ensemble Approach to X-ray Polarimetry}}\\ \textbf{A.L.Peirson}, R.W.Romani \href{https://arxiv.org/abs/2111.03047}{\emph{NeurIPS ML4PS workshop}}, 2021}
	\end{ressubsecPu}

\begin{ressubsecPu}{\textbf{\emph{Towards Optimal Signal Extraction for Imaging X-ray Polarimetry}}\\ \textbf{A.L.Peirson}, R.W.Romani. \href{https://iopscience.iop.org/article/10.3847/1538-4357/ac157d/pdf}{\emph{ApJ}, 920, 40}, 2021}
\end{ressubsecPu}

\begin{ressubsecPu}{\textbf{\emph{Deep Ensemble Analysis for Imaging X-ray Polarimetry}}\\ \textbf{A.L.Peirson}, R.W.Romani, H.L.Marshall, J.F.Steiner, L.Baldini.  \href{https://www.sciencedirect.com/science/article/pii/S0168900220311372?via\%3Dihub}{\emph{NIMA}, 986}, 2021}
\end{ressubsecPu}

\begin{ressubsecPu}{\textbf{\emph{Dank Learning: Generating Memes Using Deep Neural Networks}} [\href{https://github.com/alpv95/Dank-Learning}{Code}][\href{https://techcrunch.com/2018/06/15/dank-learning-system-autogenerates-memes/}{TechCrunch}][\href{https://thenextweb.com/artificial-intelligence/2018/06/15/stanford-researchers-taught-ai-to-make-dank-memes/}{The Next Web}]\\ \textbf{A.L.Peirson} \& E.M.Tolunay, \href{https://arxiv.org/abs/1806.04510}{1806.04510}, 2018}
\end{ressubsecPu}

\vspace{10pt}
\resitem{\textbf{\textcolor{Mycolor2}{--- Astrophysics ---}}}
\vspace{2pt}

\begin{ressubsecPu}{\textbf{\emph{X-ray Polarization of BL Lacertae in Outburst}}\\ \textbf{A.L.Peirson}, M.Negro, I.Liodakis et al. \href{https://doi.org/10.3847/2041-8213/acd242}{\emph{ApJ Letters}, 948, L25}, 2023}
\end{ressubsecPu}

\begin{ressubsecPu}{\textbf{\emph{Polarized Blazar X-rays imply Particle Acceleration in Shocks}}\\ I.Liodakis, et al. (IXPE collaboration, incl. \textbf{A.L.Peirson}) \href{https://doi.org/10.1038/s41586-022-05338-0}{\emph{Nature}, 611, 677}, 2022}
\end{ressubsecPu}

\begin{ressubsecPu}{\textbf{\emph{Polarized X-rays Constrain the Disk-Jet Geometry in the Black Hole X-ray Binary Cygnus X-1}}\\ H.Krawczynski, et al. (IXPE collaboration, incl. \textbf{A.L.Peirson}) \href{https://doi.org/10.1126/science.add5399}{\emph{Science}, 378, 650}, 2022}
\end{ressubsecPu}

\begin{ressubsecPu}{\textbf{\emph{Testing High-Energy Emission Models for Blazars with X-ray Polarimetry
}}\\ \textbf{A.L.Peirson}, I.Liodakis, R.W.Romani \href{https://iopscience.iop.org/article/10.3847/1538-4357/ac6a54}{\emph{ApJ}, 931, 59}, 2022}
\end{ressubsecPu}

\begin{ressubsecPu}{\textbf{\emph{New Tests of Millilensing in the Blazar PKS 1413+135}} [\href{https://github.com/rpoleski/MulensModel}{Code}]\\ \textbf{A.L.Peirson}, I.Liodakis, A.C.S.Readhead et al. \href{https://iopscience.iop.org/article/10.3847/1538-4357/ac469e}{\emph{ApJ}, 927, 24}, 2022}
	\end{ressubsecPu}

\begin{ressubsecPu}{\textbf{\emph{The Polarization Behavior of Relativistic Synchrotron Self-Compton Jets}} [\href{https://github.com/alpv95/SSCpol}{Code}]\\ \textbf{A.L.Peirson}, R.W.Romani. \href{https://iopscience.iop.org/article/10.3847/1538-4357/ab46b1}{\emph{ApJ}, 885, 76}, 2019}
\end{ressubsecPu}

\begin{ressubsecPu}{\textbf{\emph{Prospects for Detecting X-ray Polarization in Blazar Jets}}\\ I.Liodakis, \textbf{A.L.Peirson}, R.W.Romani. \href{https://iopscience.iop.org/article/10.3847/1538-4357/ab2719}{\emph{ApJ}, 880, 29}, 2019}
\end{ressubsecPu}

\begin{ressubsecPu}{\textbf{\emph{The Polarization Behavior of Relativistic Synchrotron Jets}}\\ \textbf{A.L.Peirson}, R.W.Romani. \href{https://iopscience.iop.org/article/10.3847/1538-4357/aad69d/meta}{\emph{ApJ}, 864, 140}, 2018}
\end{ressubsecPu}

\end{ressection}

\end{document}

